% Comparação qualitativa em pmo_52 — o item escolhido pela regra declarada na
% legenda, não por pontuação. Os dois modelos descrevem o mesmo processo; o
% DF-F1 do candidato é 0,000 porque a referência nomeia atividades como frases
% em terceira pessoa e o candidato usa forma imperativa, conforme instruído.
% Diagramas gerados por `src.evaluation.bpmn_tikz` a partir do BPMNDI real.
\begin{figure}[htb]
	\Caption{\label{fig:qualitativa} Divergência de nomenclatura sem divergência de estrutura (\textit{pmo\_52})}
	\centering

	\vspace{0.2em}
	{\small\textbf{(a) Referência de especialista}}\\[0.3em]
	\resizebox{\textwidth}{!}{% pmo_52 — gold
\begin{tikzpicture}[x=1cm, y=1cm, font=\tiny]
  \draw[-{Latex[length=1.4mm]}, gray!70] (0.43,-0.91) -- (2.28,-0.91);
  \draw[-{Latex[length=1.4mm]}, gray!70] (3.72,-0.91) -- (4.56,-0.91);
  \draw[-{Latex[length=1.4mm]}, gray!70] (6.00,-0.91) -- (6.84,-0.91);
  \draw[-{Latex[length=1.4mm]}, gray!70] (7.44,-0.91) -- (9.12,-0.91);
  \draw[-{Latex[length=1.4mm]}, gray!70] (7.44,-0.91) -- (9.12,-3.08);
  \draw[-{Latex[length=1.4mm]}, gray!70] (10.56,-0.91) -- (11.40,-0.91);
  \draw[-{Latex[length=1.4mm]}, gray!70] (10.56,-3.08) -- (11.40,-0.91);
  \draw[-{Latex[length=1.4mm]}, gray!70] (12.00,-0.91) -- (13.68,-0.91);
  \node[rounded corners=2pt, draw, thin, fill=white, align=center, text width=1.66cm, minimum height=0.96cm, inner sep=1pt] at (3.00,-0.91) {MPON sents the\\dismissal to the\\MPOO};
  \node[rounded corners=2pt, draw, thin, fill=white, align=center, text width=1.66cm, minimum height=0.96cm, inner sep=1pt] at (5.28,-0.91) {MPOO reviews the\\dismissal};
  \node[rounded corners=2pt, draw, thin, fill=white, align=center, text width=1.66cm, minimum height=0.96cm, inner sep=1pt] at (9.84,-0.91) {MPOO opposes the\\dismissal};
  \node[rounded corners=2pt, draw, thin, fill=white, align=center, text width=1.66cm, minimum height=0.96cm, inner sep=1pt] at (9.84,-3.08) {MPOO confirmes the\\dismissal};
  \node[diamond, draw, thin, minimum size=0.52cm, inner sep=0pt, fill=white] at (7.14,-0.91) {$\times$};
  \node[diamond, draw, thin, minimum size=0.52cm, inner sep=0pt, fill=white] at (11.70,-0.91) {$\times$};
  \node[circle, draw, thin, minimum size=0.42cm, inner sep=0pt, fill=white] at (0.22,-0.91) {};
  \node[circle, draw, very thick, minimum size=0.42cm, inner sep=0pt, fill=white] at (13.90,-0.91) {};
\end{tikzpicture}
}

	\vspace{1.4em}
	{\small\textbf{(b) Braço A4 --- modelo de 7B com ajuste fino \quad DF-F1 $= 0{,}000$}}\\[0.3em]
	\resizebox{\textwidth}{!}{% pmo_52 — A4
\begin{tikzpicture}[x=1cm, y=1cm, font=\tiny]
  \draw[-{Latex[length=1.4mm]}, gray!70] (0.43,-0.91) -- (2.28,-0.91);
  \draw[-{Latex[length=1.4mm]}, gray!70] (3.72,-0.91) -- (4.56,-5.81);
  \draw[-{Latex[length=1.4mm]}, gray!70] (7.44,-5.81) -- (9.12,-5.81);
  \draw[-{Latex[length=1.4mm]}, gray!70] (10.56,-5.81) -- (11.40,-5.81);
  \draw[-{Latex[length=1.4mm]}, gray!70] (11.83,-5.81) -- (13.68,-5.81);
  \draw[-{Latex[length=1.4mm]}, gray!70] (7.44,-5.81) -- (9.12,-7.98);
  \draw[-{Latex[length=1.4mm]}, gray!70] (10.56,-7.98) -- (11.40,-7.98);
  \draw[-{Latex[length=1.4mm]}, gray!70] (11.83,-7.98) -- (13.68,-5.81);
  \draw[-{Latex[length=1.4mm]}, gray!70] (6.00,-5.81) -- (6.84,-5.81);
  \node[circle, draw, thin, minimum size=0.42cm, inner sep=0pt, fill=white] at (0.22,-0.91) {};
  \node[anchor=north, font=\tiny] at (0.22,-1.19) {Start};
  \node[rounded corners=2pt, draw, thin, fill=white, align=center, text width=1.66cm, minimum height=0.96cm, inner sep=1pt] at (3.00,-0.91) {Send Dismissal};
  \node[rounded corners=2pt, draw, thin, fill=white, align=center, text width=1.66cm, minimum height=0.96cm, inner sep=1pt] at (5.28,-5.81) {Review Dismissal};
  \node[diamond, draw, thin, minimum size=0.52cm, inner sep=0pt, fill=white] at (7.14,-5.81) {$\times$};
  \node[anchor=south, font=\tiny] at (7.14,-5.49) {Does the MPOO oppose};
  \node[diamond, draw, thin, minimum size=0.52cm, inner sep=0pt, fill=white] at (13.98,-5.81) {$\times$};
  \node[anchor=south, font=\tiny] at (13.98,-5.49) {Does the MPOO oppose};
  \node[rounded corners=2pt, draw, thin, fill=white, align=center, text width=1.66cm, minimum height=0.96cm, inner sep=1pt] at (9.84,-5.81) {Oppose Dismissal};
  \node[circle, draw, very thick, minimum size=0.42cm, inner sep=0pt, fill=white] at (11.62,-5.81) {};
  \node[anchor=north, font=\tiny] at (11.62,-6.09) {End Opposed};
  \node[rounded corners=2pt, draw, thin, fill=white, align=center, text width=1.66cm, minimum height=0.96cm, inner sep=1pt] at (9.84,-7.98) {Confirm Dismissal};
  \node[circle, draw, very thick, minimum size=0.42cm, inner sep=0pt, fill=white] at (11.62,-7.98) {};
  \node[anchor=north, font=\tiny] at (11.62,-8.26) {End Confirmed};
\end{tikzpicture}
}

	\Fonte{Elaborado pelo autor}
	\Nota{Item selecionado pela regra, declarada antes da inspeção, de \textbf{menor
	número de nós na referência entre os itens em que o braço A4 produziu saída
	válida} --- critério de legibilidade, independente de pontuação. Os diagramas
	usam a topologia e o leiaute automáticos do próprio \textit{pipeline}; apenas
	o espaçamento vertical foi ampliado, porque os rótulos da referência ocupam
	quatro a cinco linhas e estouram a altura reservada ao nó pelo leiaute.
	A referência escreve \textit{MPOO confirmes the dismissal}; o candidato,
	\textit{Confirm Dismissal}. A sobreposição de \textit{tokens} normalizados fica
	abaixo do limiar de casamento, e a aresta correspondente é contabilizada como
	erro estrutural.}
\end{figure}
